This paper presented the design, implementation, and field deployment of a low-cost, low-power LoRa wireless sensor network for monitoring environmental conditions beneath Icelandic turf roofs. The system integrates distributed sensor nodes with local SD-card storage, 433~MHz LoRa communication with mesh relay functionality, a base station, and the ThingSpeak cloud platform. Four nodes were deployed beneath four turf roofs at Keldur and delivered valid records over periods of approximately 51 to 130 days.

Under the defined operating cycle, the average node current was approximately 3.08~mA, giving an idealized lifetime of about 142 days from 10,500~mAh. This is a calculated estimate that excludes low-temperature capacity loss, and the logged operating periods of the single-cell preliminary tests (approximately 11--27 days) are not full battery lifetimes. Single-hop LoRa communication was observed at several locations within about 3~km, with occasional packet loss in obstructed areas. Mesh relaying was observed only qualitatively on the bench, so no multi-hop range is reported. The soil moisture calibration produced a second-order polynomial with an in-sample RMSE of approximately 3.01 percentage points over the 0--40\% gravimetric range, and ThingSpeak uploads took approximately 1--2~s.

Future work should extend the duration of field deployments, in particular over winter, and improve enclosure and sensor protection against moisture, freezing, and snow. It should add remote system-health monitoring and base-station logging so that missing LoRa transmissions can be distinguished from cloud-upload interruptions. It should also recalibrate the soil moisture sensor with soil from the deployment sites against a separate validation dataset and carry out controlled multi-hop tests with retained logs.
